\chapter{Частые ошибки}\label{app:oshibki}

Сообщения ниже настоящие: их получили, запустив файлы из \kod{solutions/bad_examples}. Номер в скобках — строка файла, а не шоу.

\section*{Ошибки запуска кода}

\begin{center}
\small
\begin{longtable}{@{}p{6.6cm}p{7.4cm}@{}}
\toprule
Сообщение & Что случилось и что делать \\
\midrule
\endhead
\kod{Синтаксическая ошибка в строке 6: '(' was never closed} & Не закрыта скобка (\kod{bad_syntax.py}). Посмотрите на указанную строку и строку выше. \\
\kod{NameError: name 'e' is not defined (строка 7)} & Написано имя, которое не создано, обычно опечатка: \kod{e.land()} вместо \kod{d.land()}. \\
\kod{ValueError: Неизвестный цвет 'pink'. Доступны: black, red, green, yellow, blue, purple, cyan, white (строка 6)} & Цвета \kod{pink} нет. Возьмите имя из списка, номер 0–7 или тройку \kod{(r, g, b)}. \\
\kod{ValueError: land_on: сначала нужен взлёт (строка 6)} & \kod{land_on} берёт высоту у последнего \kod{takeoff}, \kod{set_z} или \kod{code_z}. Поставьте их раньше. \\
\kod{В файле нет переменной show = Show('имя')} & Нет строки \kod{show = Show('имя')} или переменная названа иначе. \\
\kod{Шоу содержит больше 10000 шагов. Проверьте циклы в коде} & Цикл без конца или с огромным числом повторов (\kod{bad_loop.py}). \\
\kod{Шоу выполняется дольше 5 с. Проверьте циклы в коде} & Код шоу работает слишком долго; ищите бесконечный цикл. \\
\kod{Время шага не может быть отрицательным} & Отрицательное число секунд в команде. \\
\kod{Список клеток пуст} & В \kod{code_s} передан пустой список. \\
\kod{Номер цвета должен быть от 0 до 7} и \kod{Компоненты цвета должны быть от 0 до 255} & Номер цвета или канал вне допустимого диапазона. \\
\bottomrule
\end{longtable}
\end{center}

\section*{Сообщения проверки}

\begin{center}
\small
\begin{longtable}{@{}p{6.6cm}p{7.4cm}@{}}
\toprule
Сообщение & Что делать \\
\midrule
\endhead
\kod{Дрон 1: клетка 144 вне коврика 0-143} & Номер клетки может быть от 0 до 143. \\
\kod{Дрон 1: высота 180 см больше лимита 150 см, полёт будет отклонён} & Уменьшите высоту или поднимите лимит в окошке \knopka{Лимит высоты} (не больше 200\cm), потом снова \knopka{Проверить}. \\
\kod{Дрон 1: переход на клетку 16 за 1 с (120 см, 120 см/с) быстрее порога 100 см/с: дрон может перелететь клетку} & Добавьте секунд: расстояние / секунды не больше 100. \\
\kod{Дроны 1 и 2 сближаются до 2 см на 12.3 с (порог 60 см)} & Разведите дронов по высоте, по времени или по маршруту. \\
\kod{Зависание 90 с больше лимита 60 с} & Ожидание на месте длиннее 60\,с (отказ при загрузке). Разбейте его или проверьте число. \\
\kod{Дрон N: полёт выходит за рабочую область (…)} & Дрон вылетел за выбранную область. Смените область, нажмите \knopka{Перенести шоу} или поправьте клетки. \\
\kod{Шаги circle ещё не летали на живом дроне} & Не ошибка, а предупреждение: на живом дроне эти шаги не проверены, время круга будет другим. Первый полёт делайте одним дроном. \\
\kod{загрузка в дрон недоступна} & В программе дрона есть шаг, который загрузить нельзя (\kod{circle_z}). \\
\kod{Перед загрузкой поставьте дрон на его первую клетку} & Сведение: поставьте дрон на первую клетку шоу и сделайте взлёт на 50\cm. \\
\bottomrule
\end{longtable}
\end{center}

\section*{Ошибки, которые проверка не видит}

\begin{itemize}
  \item \textbf{Забыли ожидание.} Дрон 2 должен начать «после» дрона 1, но у них одновременный старт; вы не вставили \kod{set_z} на той же высоте.
  \item \textbf{Перепутали клетки.} Номер клетки не столбец и не строка, а столбец + 12 $\times$ строка. Клетка 16 — это строка 1, а не 6-я клетка.
  \item \textbf{Не тот дрон шоу.} Живой дрон стоит не на первой клетке нужного дрона: загрузка откажет или дрон полетит не туда.
  \item \textbf{Цвет без ожидания.} \kod{rgb} занимает 0\,с: если следом идёт ещё один \kod{rgb}, первый цвет не успеет показаться.
  \item \textbf{Два дрона на одной клетке.} В симуляторе дроны стартуют под своей первой клеткой; два дрона на одной клетке дадут замечание о сближении.
\end{itemize}
